\section{A Linear System Perspective on the Canonical PSO}
\label{sec:decomposition}

\cite{4424687} imports a linear system perspective to the canonical PSO. 
There are two components that influence the systematic dynamic, which are
\begin{itemize}
\item how the global best and personal best are updated;
\item the parameters in the update rule.
\end{itemize}
In designing a PSO algorithm for the properties in different problem, different strategies are applied to update the global best and the personal best.
Relatively, the update rule is rarely modified.

Without loss of generality, we look at only a one-dimension particle to simplify the equations.
\begin{subequations}
\label{eq:pso1_alg}
\begin{equation}
\label{eq:up_vel1}
\begin{aligned}
v(k+1) & = \chi [ v(k) 
 + \phi^{P} u^{P}(k) (x^{P}(k) - x(k)) + \phi^{G} u^{G}(k) (x^{G}(k) - x(k)) ] \\
& = \chi v(k) + ( - \chi \phi^{P} u^{P}(k) - \chi \phi^{G} u^{G}(k) ) x(k) + \phi^{P} u^{P}(k) x^{P}(k) + \phi^{G} u^{G}(k) x^{G}(k),
\end{aligned}
\end{equation}
\begin{equation}
\label{eq:up_pos1}
\begin{aligned}
x(k+1) & = x(k) + v(k+1) \\
& = \chi v(k) + (1 - \chi \phi^{P} u^{P}(k) - \chi \phi^{G} u^{G}(k) ) x(k) + \phi^{P} u^{P}(k) x^{P}(k) + \phi^{G} u^{G}(k) x^{G}(k) .
\end{aligned}
\end{equation}
\end{subequations}

We can organize the equations \eqref{eq:pso1_alg} into a linear form.

\begin{equation}
\label{eq:pso_up_linalg}
\begin{aligned}
\begin{bmatrix}
v(k+1)
\\ 
x(k+1)
\end{bmatrix}
= &
\begin{bmatrix}
\chi & - \chi (\phi^{G} u^{G}(k) + \phi^{P} u^{P}(k) )
\\ 
\chi & 1 - \chi (\phi^{G} u^{G}(k) + \phi^{P} u^{P}(k) )
\end{bmatrix}
\begin{bmatrix}
v(k)
\\ 
x(k)
\end{bmatrix}
\\ & +  
\begin{bmatrix}
\chi \phi^{G} u^{G}(k) & \chi \phi^{P} u^{P}(k)
\\ 
\chi \phi^{G} u^{G}(k) & \chi \phi^{P} u^{P}(k)
\end{bmatrix}
\begin{bmatrix}
x^{G}(k)
\\ 
x^{P}(k)
\end{bmatrix}.
\end{aligned}
\end{equation}

Define
\begin{equation}
\label{eq:func_a}
a(k) = \chi \phi^{G} u^{G}(k)
\end{equation}
and
\begin{equation}
\label{eq:func_b}
b(k) = \chi \phi^{P} u^{P}(k),
\end{equation}
we have a simplified form of equation \eqref{eq:pso_up_linalg} as
\begin{equation}
\label{eq:pso_up_linalg_simp}
\begin{aligned}
\begin{bmatrix}
v(k+1)
\\ 
x(k+1)
\end{bmatrix}
& =
\begin{bmatrix}
\chi & - a(k) - b(k)
\\ 
\chi & 1 - a(k) - b(k)
\end{bmatrix}
\begin{bmatrix}
v(k)
\\ 
x(k)
\end{bmatrix}
\\
& + 
\begin{bmatrix}
a(k) & b(k)
\\ 
a(k) & b(k)
\end{bmatrix}
\begin{bmatrix}
x^{G}(k)
\\ 
x^{P}(k)
\end{bmatrix},
\end{aligned}
\end{equation}
in which
$ a(k) \in [0, \chi \phi^{G} ] $ and $ b(k) \in [0, \chi \phi^{P}] $.

\begin{figure}
\centering
\includegraphics[width=0.6\linewidth]{sys_flow.pdf}
\caption{A system structure of the PSO}
\label{fig:sys_flow}
\end{figure}

We can decompose the model of a single particle into two components, which are illustrated in Figure \ref{fig:sys_flow}.
Two components forms a cascade system structure.
An \emph{optimization strategy sub-system} generates $ x^{G}_{i}(k) $ and $ x^{P}_{i}(k) $ based on last position $ x_{i}(k) $ and the swarm information from the topology interaction.
$ x^{G}_{i}(k) $ and $ x^{P}_{i}(k) $ are used as inputs to the \emph{rule update sub-system}.

The convergence problem are answered by
\begin{itemize}
\item whether the \emph{rule update sub-system} is input-to-state stable;
\item how the output of the \emph{optimization strategy sub-system} is like.
\end{itemize}

If the \emph{rule update sub-system} is input-to-state stable, the convergence of $ x_{i}(k) $ depends on the boundaries of $ x^{G}_{i}(k) $ and $ x^{P}_{i}(k) $.
In the next section, we will evaluate the input-to-state stability of the rule update sub-system.

